%
% teil2.tex -- Polardarstellung und Drehung
%
\section{Polardarstellung und Drehung
\label{komplex:section:polar}}
\kopfrechts{Polardarstellung und Drehung}
Die kartesische Form $z = a + bi$ ist für die Addition optimal,
für die Multiplikation aber unhandlich.
Die Polardarstellung schliesst diese Lücke
und macht zugleich die geometrische Deutung der Multiplikation als
Drehstreckung sichtbar.
In den folgenden Unterabschnitten werden Polarform, Multiplikation als
Rotation und schliesslich der Drehfaktor entwickelt,
mit dem alle späteren Konstruktionen ausgeführt werden.

%
% Polarform einer komplexen Zahl
%
\subsection{Polarform einer komplexen Zahl
\label{komplex:subsection:polarform}}
Jeder Punkt $z \neq 0$ der Ebene ist eindeutig durch seinen Abstand
$r = |z|$ vom Ursprung
und den Winkel $\varphi$, den der Pfeil zum Punkt mit der positiven reellen
Achse einschliesst, festgelegt.
Mit
\begin{equation}
z
=
r \cdot ( \cos(\varphi) + i \sin(\varphi) )
\label{komplex:eq:polarform}
\end{equation}
erhält man die \emph{Polarform}.
\index{Polarform}%
Der Winkel $\varphi$ heisst \emph{Argument} und wird mit $\arg(z)$
\index{Argument}%
bezeichnet;
er ist nur bis auf ganzzahlige Vielfache von $2\pi$ eindeutig.
Üblicherweise wählt man den Hauptwert in $(-\pi, \pi]$.
Die Umrechnung zwischen kartesischer und polarer Form lautet
\begin{equation}
x = r \cdot \cos(\varphi) ,
\qquad
y = r \cdot \sin(\varphi) ,
\qquad
r = \!\sqrt{x^2 + y^2} .
\label{komplex:eq:umrechnung}
\end{equation}
Das Vorzeichen des Winkels gibt den Drehsinn an:
$\varphi > 0$ entspricht einer Drehung gegen den Uhrzeigersinn,
$\varphi < 0$ einer Drehung im Uhrzeigersinn.

%
% Multiplikation als Rotation
%
\subsection{Multiplikation als Rotation
\label{komplex:subsection:rotation}}
Aus den Additionstheoremen für Sinus und Kosinus folgt für zwei Zahlen in
Polarform $z_1 = r_1 ( \cos(\varphi_1) + i \sin(\varphi_1) )$
und $z_2 = r_2 ( \cos(\varphi_2) + i \sin(\varphi_2) )$ die Produktformel
\begin{equation}
z_1 \cdot z_2
=
r_1 r_2 \cdot \bigl( \cos(\varphi_1 + \varphi_2)
                   + i \sin(\varphi_1 + \varphi_2) \bigr) .
\label{komplex:eq:produkt-polar}
\end{equation}
Diese Identität ist die zentrale Aussage des gesamten Abschnitts:
Bei der Multiplikation komplexer Zahlen multiplizieren sich die Beträge,
und die Argumente addieren sich.
Die Multiplikation mit einer festen Zahl
$w = r \cdot ( \cos(\varphi) + i \sin(\varphi) )$
wirkt damit auf die Ebene als \emph{Drehstreckung} ---
eine Streckung um den Faktor $r$ und eine gleichzeitige Drehung um den
Winkel $\varphi$, beides um den Ursprung.

Insbesondere ist die Multiplikation mit
$\cos(\varphi) + i \sin(\varphi)$
eine reine Drehung um den Ursprung um den Winkel $\varphi$.
Diese Beobachtung verwandelt jede Drehung in der Ebene in eine algebraische
Operation, was in den späteren Konstruktionen zur Schlüsselrechnung wird.

%
% Der Drehfaktor
%
\subsection{Der Drehfaktor
\label{komplex:subsection:drehfaktor}}
Im Folgenden bezeichnen wir mit $d$ den \emph{Drehfaktor}
\index{Drehfaktor}%
\begin{equation}
d
=
\cos(\varphi) + i \sin(\varphi) ,
\label{komplex:eq:drehfaktor}
\end{equation}
also die komplexe Zahl vom Betrag $1$, deren Multiplikation einer Drehung
um den Winkel $\varphi$ entspricht.
Eine Drehung um einen beliebigen Punkt $z_0$ lässt sich aus drei Schritten
zusammensetzen:
Verschiebung des Drehzentrums in den Ursprung,
Drehung um den Ursprung
und Verschiebung zurück.
In Formeln gilt
\begin{equation}
D(z)
=
z_0 + d \cdot (z - z_0) .
\label{komplex:eq:drehung-allg}
\end{equation}
Diese Formel wird im Folgenden so häufig verwendet,
dass sie als eigentliches Arbeitspferd der Arbeit gelten kann.
Sie ist die direkte algebraische Übersetzung des Konzepts ``Drehung um
einen Punkt'' und erlaubt es, Konstruktionen, die mit Zirkel und Lineal
sehr aufwändig wären, in wenigen Zeilen zu rechnen.
